% from ega4-2-fr.tex
\oldpage[IV-2]{59}

\begin{proposition}[4.3.9]
\label{IV.2.4.3.9-fr}
Si $k$ est algébriquement clos, alors, pour deux extensions quelconques $K$,
$L$ de $k$, $L\otimes_k K$ est intègre, et réciproquement.
\end{proposition}

\begin{proof}
C'est une conséquence de (4.3.3) et (4.3.6), un corps parfait et séparablement
clos étant algébriquement clos.
\end{proof}

\subsection{Préschémas irréductibles et préschémas connexes sur un corps
algébriquement clos}
\label{subsection:IV.2.4.4-fr}

\begin{env}[4.4.1]
\label{IV.2.4.4.1-fr}
Soient $k$ un corps, $X$ un $k$-préschéma, $K$ une extension de $k$. Comme le
morphisme $\operatorname{Spec}(K)\to\operatorname{Spec}(k)$ est fidèlement
plat, quasi-compact, et universellement ouvert (2.4.9), il en est de même du
morphisme projection $p:X\otimes_k K\to X$ (2.2.13). Il résulte donc de
(2.3.5) que toute composante irréductible de $X\otimes_k K$ \emph{domine} une
composante irréductible de $X$, et (comme $p$ est surjectif) toute composante
irréductible de $X$ est dominée par une composante irréductible de
$X\otimes_k K$~; on déduit ainsi de $p$ une application \emph{surjective} de
l'ensemble des composantes irréductibles de $X\otimes_k K$ dans l'ensemble
des composantes irréductibles de $X$. De même, il résulte du fait que $p$ est
continu et surjectif que toute composante connexe de $X$ est image par $p$
d'une réunion de composantes connexes de $X\otimes_k K$, d'où une application
\emph{surjective} de l'ensemble des composantes connexes de $X\otimes_k K$
dans l'ensemble des composantes connexes de $X$. On en conclut que toute
composante irréductible (resp. connexe) $Z'$ de $X\otimes_k K$ est contenue
dans un \emph{unique} ensemble de la forme $p^{-1}(Z)$, où $Z$ est une
composante irréductible (resp. connexe) de $X$~; pour tout sous-préschéma de
$X$ ayant $Z$ pour espace sous-jacent (et encore noté $Z$), $Z'$ est une
composante irréductible (resp. connexe) de $Z\otimes_k K$. En particulier, si
$X\otimes_k K$ a un nombre fini $n$ (resp. $n'$) de composantes irréductibles
(resp. connexes) (ce qui sera le cas lorsque $X$ est \emph{de type fini sur
$k$}, car alors $X\otimes_k K$ est de type fini sur $K$, donc noethérien), le
nombre de composantes irréductibles (resp. connexes) de $X$ est $\leqslant n$
(resp. $\leqslant n'$)~; dire qu'il est égal à $n$ (resp. $n'$) signifie que
pour tout sous-préschéma réduit $Z$ de $X$ ayant pour espace sous-jacent une
composante irréductible (resp. connexe) de $X$, $Z\otimes_k K$ est
irréductible (resp. connexe) (I, 5.1.8)~; le nombre de composantes
irréductibles (resp. connexes) de $X\otimes_k K$ est donc indépendant de $K$
si et seulement si pour toute composante irréductible (resp. connexe) $Z$ de
$X$, $Z\otimes_k K$ est irréductible (resp. connexe) pour \emph{toute}
extension $K$ de $k$.

En particulier, si $X\otimes_k K$ est irréductible (resp. connexe), il en est
de même de $X$ (ce qui résulte déjà du fait que $p$ est surjectif). De même, si
$X\otimes_k K$ est \emph{réduit} (resp. \emph{intègre}), il en est de même de
$X$ puisque $p$ est fidèlement plat (2.1.13).

Dans ce numéro et le suivant, nous allons examiner de plus près ce qu'on peut
dire des composantes irréductibles (resp. connexes) de $X\otimes_k K$ lorsque
$K$ varie~; ce qui précède nous montre déjà que le nombre de ces composantes
est fonction \emph{croissante} de $K$.
\end{env}

\begin{lemma}[4.4.2]
\EGAReaderTarget{ega4.c04.lem.4-4-2}{IV.2.4.4.2-fr}
Soient $X'$, $X$ deux espaces topologiques, $f:X'\to X$ une application
continue. On suppose que les conditions suivantes sont vérifiées~:
\begin{enumerate}
\item[(i)] $f$ est surjective ouverte (resp. telle que l'espace topologique
$X$ s'identifie canoniquement au quotient de $X'$ par la relation d'équivalence
définie par $f$).

